\section{Planned Library}
ReMix will expose the four estimators through an interface accepting a mixture
and an objective. Diagonal Gaussian components will have trainable logits, means,
and positive scales. Each estimator will return an objective estimate for logging
and a scalar surrogate for automatic differentiation, with optional STL for
the ELBO.

Table~\ref{tab:estimators} summarizes the requirements. Each mixture density
evaluation includes all $K$ components; stratification repeats this work at every
sample point. Shared operations will use log-sum-exp for mixture log densities,
normalized component log densities for responsibilities, and triangular solves
for implicit derivatives.

\begin{table}[htbp]
 \centering
 \caption{Estimator requirements. Stratification uses one sample per component.}
 \label{tab:estimators}
 \small
 \begin{tabular}{@{}lcl@{}}
  \toprule
  Estimator & Objective evaluations & Required operations \\
  \midrule
  Score-function & $1$ & Mixture log-density derivatives \\
  Stratified & $K$ & Component sampling derivatives \\
  Implicit & $1$ & Conditional CDFs and a Jacobian solve \\
  Post-stratified & $1$ & Inverse component maps and responsibilities \\
  \bottomrule
 \end{tabular}
\end{table}

We will compare gradient variance, ELBO progress, runtime, and memory use across
component counts, overlap levels, and event dimensions. Progress per model
evaluation and per unit of time will measure practical optimization efficiency.
